% main_si.tex — Supporting Information. Compile with lualatex (Lean listings need Unicode).
% Tables ST1-ST3 are generated from the data by make_si_tables.py; other numbers cite ledger IDs in % comments.
\documentclass[11pt]{article}
\usepackage[a4paper,margin=2.2cm]{geometry}
\usepackage{fontspec}
\setmonofont{DejaVu Sans Mono}[Scale=0.82]
\usepackage{amsmath,amssymb}
\usepackage{booktabs,longtable,seqsplit}
\usepackage{fvextra}
\usepackage{graphicx}
\usepackage[hidelinks]{hyperref}
% notation.tex — macros used across the manuscript (hard rule 1)
\usepackage{siunitx}
\sisetup{separate-uncertainty=true}
\newcommand{\Eks}{E_{\mathrm{KS}}}          % Kohn--Sham HOMO--LUMO gap
\newcommand{\Eopt}{E_{\mathrm{opt}}}        % optical gap (experimental or calibrated)
\newcommand{\Esone}{E_{\mathrm{S_1}}}       % lowest singlet excitation energy
\newcommand{\Voc}{V_{\mathrm{OC}}}
\newcommand{\Jsc}{J_{\mathrm{SC}}}
\newcommand{\FF}{\mathit{FF}}
\newcommand{\PCE}{\mathit{PCE}}
\newcommand{\MAE}{\mathrm{MAE}}
\newcommand{\sTDA}{sTDA}
\newcommand{\sTDAxTB}{sTDA-xTB}
\usepackage{tikz}

\renewcommand{\thesection}{S\arabic{section}}
\renewcommand{\thetable}{S\arabic{table}}
\renewcommand{\thefigure}{S\arabic{figure}}

\title{Supporting Information\\[4pt]\large Class-Dependent Errors of Cheap Band-Gap Observables\\ Undermine
High-Throughput Screening of Organic Photovoltaic Absorbers}
\author{Patrick Sorrel Mvoto Kongo, Steve Cabrel Teguia Kouam,\\ Jean-Pierre Tchapet Njafa, Serge Guy Nana Engo}
\date{}

\begin{document}
\maketitle
\tableofcontents

\section{Validity rules of gate G0}
Table~\ref{tab:si-validity} lists how many PubChemQC records each validity rule removes. Two rules dominate: 6\,323
records fail the ring or aromaticity rules, and 838 records contain more than one fragment. The second group matters
beyond its size, because neither gap observable describes a single absorber for a salt or a mixture
(Section~\ref{sec:salts}).
% L-2026-10-17-08
\begin{table}[h]
\centering
\caption{Validity rules of gate G0, applied in the order listed; each record is counted under the first rule it fails.
10\,215 of 17\,458 records pass.}
\label{tab:si-validity}
\begin{tabular}{llr}
\toprule
Rule & Record removed when the structure has & Records\\
\midrule
R1 & more than one fragment (salts\, mixtures\, co-crystals) & 838 \\
R2 & contains a metal & 2 \\
R3 & no ring & 4\,034 \\
R4 & no aromatic ring & 2\,289 \\
R5 & fewer than six heavy atoms & 9 \\
R6 & fewer than six conjugated atoms & 70 \\
-- & SMILES not parsable by RDKit & 1 \\
\bottomrule
\end{tabular}
\end{table}


\section{Salts, not halogens, at the window edges}\label{sec:salts}
Halogenated records are more frequent near the absorber window than in the database as a whole, which could be read as
a chemical trend. Table~\ref{tab:si-salts} shows that the excess follows the salts. The database contains no bromine or
iodine; its chlorine is bound covalently, present as an isolated fragment (chloride or HCl), or both. Covalently bound
chlorine occurs in 2 of the 33 records of the Kohn--Sham window (6.1\,\%, both of them hydrochloride salts) and in
the sTDA-xTB window at about its database frequency (13.4 against 14.7\,\%). The excess is carried by multi-fragment
records: they make up 48.5\,\% of the Kohn--Sham window and 84.4\,\% of the records whose lowest sTDA-xTB singlet
lies below \qty{1}{eV}, consistent with spurious charge-transfer states between the separated ions. After Holm correction for the four enrichment tests performed, the halogen excess in the
Kohn--Sham window is not significant ($p = 0.11$), while the excess among the collapsed records is ($p = 7 \times
10^{-127}$). Removing multi-fragment records (rule R1) therefore halves the Kohn--Sham window and removes most
collapse cases.
% L-2026-10-17-05 (outputs/extract/halogen-frontier.txt); Holm: L-2026-10-18-12 (stat-hardening.txt)
\begin{table}[h]
\centering
\caption{Composition of the 15\,952 records with both observables and of three subsets. Percentages of each set.
A record can count both as containing an isolated Cl fragment and as containing covalent Cl. Odds ratio and
Holm-adjusted Fisher $p$ for halogenated records in the subset against the rest of the database.}
\label{tab:si-salts}
\resizebox{\textwidth}{!}{\begin{tabular}{lrrrrrrr}
\toprule
Set & $n$ & Halogenated & Multi-fragment & Isolated Cl fragment & Covalent Cl & Odds ratio & $p_\mathrm{Holm}$\\
\midrule
All records & 15\,952 & 21.7 & 4.5 & 2.9 & 14.7 & -- & -- \\
$\Eks$ in \qtyrange{1.1}{2.2}{eV} & 33 & 36.4 & 48.5 & 36.4 & 6.1 & 2.06 & 0.11 \\
$\Esone$ in \qtyrange{1.1}{2.2}{eV} & 559 & 22.2 & 10.0 & 5.9 & 13.4 & 1.03 & 0.79 \\
$\Esone < \qty{1}{eV}$ & 742 & 61.5 & 84.4 & 53.9 & 12.0 & 6.46 & $7 \times 10^{-127}$ \\
\bottomrule
\end{tabular}}
\end{table}

\section{Basis set and geometry of the Kohn--Sham gap}\label{sec:basis}
The calibration set uses B3LYP/def2-SVP gaps on DFT geometries, the database stores B3LYP/6-31G* gaps on PM6 geometries,
and our recomputations used force-field (UFF) geometries. Table~\ref{tab:si-acenes} compares these levels for the
acenes. For pentacene at a fixed UFF geometry, def2-SVP and 6-31G* give gaps of 1.884 and \qty{1.896}{eV}: the basis
set moves the gap by \qty{0.01}{eV}, far less than the offsets discussed in the main text. Geometry moves it more:
re-optimising pentacene at B3LYP/6-31G* raises its gap to \qty{2.209}{eV}, \qty{0.31}{eV} above the UFF value at the same
level. Across the four acenes, the UFF/def2-SVP gaps lie 0.37--\qty{0.51}{eV} below the database values, a difference
that combines geometry and basis. On DFT or PM6
geometries the Kohn--Sham offset to the measured $L_a$ energy ranges from about zero (pentacene, own optimisation) to
$+0.46$~eV (naphthalene, database value), against $+0.29$~eV for the donor--acceptor molecules of HOPV15, a smaller
class contrast than for sTDA-xTB, which reproduces the acenes to \qty{0.12}{eV} on UFF geometries but overestimates the
donor--acceptor gaps by \qty{1.0}{eV}. No sTDA-xTB calculation on DFT geometries of the acenes was made, so the size of
the sTDA-xTB contrast after a geometry correction is not quantified here.
% L-2026-10-18-15 (outputs/extract/acene-def2svp.txt); geometry wording per L-2026-10-18-20
\begin{table}[h]
\centering
\caption{Kohn--Sham gaps of four acenes [eV]. Experimental values: midpoints of the measured lowest bright ($L_a$)
excitation energies (Malloci et al., 2007). Offset: Kohn--Sham gap minus experiment.}
\label{tab:si-acenes}
\resizebox{\textwidth}{!}{\begin{tabular}{lrrrrr}
\toprule
Acene & def2-SVP, UFF & 6-31G*, PM6 (PubChemQC) & Exp.\ $L_a$ & Offset, def2-SVP & Offset, PubChemQC\\
\midrule
Naphthalene & 4.537 & 4.909 & 4.445 & $+0.092$ & $+0.464$ \\
Anthracene & 3.247 & 3.709 & 3.360 & $-0.113$ & $+0.349$ \\
Tetracene & 2.431 & 2.931 & 2.660 & $-0.229$ & $+0.271$ \\
Pentacene & 1.884 & 2.397 & 2.215 & $-0.331$ & $+0.182$ \\
\bottomrule
\end{tabular}}
\end{table}

\section{Simplified TDA on a B3LYP ground state}\label{sec:std2}
\textit{Interface.} std2 reads the ground state from a Molden file. Two details decide whether the result is correct:
the \texttt{[Atoms]} block must be written in \AA ngstr\"om (\texttt{[Atoms] Angs}), because the program misreads
coordinates labelled in atomic units and then reports a singular overlap matrix; and the molecular-orbital coefficients
must refer to Cartesian Gaussian functions, so the PySCF calculation was run with \texttt{cart=True}. With these settings
std2 gave a pentacene singlet of \qty{2.094}{eV} (sTDA) against \qty{2.203}{eV} from a full TDA-B3LYP/6-31G* calculation
in PySCF on the same geometry, a difference of \qty{0.11}{eV}. Excitation energies were re-extracted from every std2 output with a parser that requires the state-table
header; all 36 molecules parsed.
% L-2026-10-18-14; outputs/literature/std2-assessment.md

\textit{Result.} Table~\ref{tab:si-std2} compares the observables on the 36 HOPV15 molecules with DFT geometries.
Moving sTDA onto a B3LYP ground state removes the \qty{0.76}{eV} overestimate of sTDA-xTB, which points to the
tight-binding ground state as the main source of that bias. In this set, sTDA@B3LYP is then the most accurate gap
observable \emph{without} calibration, with 28 of 31 molecules recalled by the raw window. Its molecule-to-molecule
scatter is larger than that of $\Eks$, however, and after a linear calibration the stored Kohn--Sham gap is better by
\qty{0.06}{eV} (paired bootstrap 95\,\% interval 0.031--\qty{0.091}{eV}, 34 complete cases). The exact-integral variant
XsTDA brings no gain here. For a screen of this class, a calibrated $\Eks$, already stored in the database, is the more
accurate choice; sTDA@B3LYP is useful as an uncalibrated cross-check.
% L-2026-10-18-14 (paired CI); table: make_si_tables.py
\begin{table}[h]
\centering
\caption{Gap observables computed on the DFT geometries distributed with HOPV15, against the experimental optical gap
($n = 36$ donor--acceptor small molecules). SD: standard deviation of the signed error. Recall: molecules whose
experimental and computed gaps both lie in the \qtyrange{1.1}{2.2}{eV} window. LOO MAE: leave-one-out error of a linear
calibration against experiment.}
\label{tab:si-std2}
\begin{tabular}{lrrrrrr}
\toprule
Observable & Bias [eV] & MAE [eV] & SD [eV] & $r$ & Recall & LOO MAE [eV]\\
\midrule
$\Eks$, B3LYP/def2-SVP & +0.29 & 0.29 & 0.16 & 0.92 & 23/31 & 0.083 \\
sTDA@B3LYP/6-31G* & -0.02 & 0.21 & 0.26 & 0.75 & 28/31 & 0.141 \\
XsTDA@B3LYP/6-31G* & +0.24 & 0.25 & 0.29 & 0.69 & 23/31 & 0.157 \\
sTDA-xTB & +0.76 & 0.76 & 0.12 & 0.89 & 0/31 & 0.103 \\
\bottomrule
\end{tabular}
\end{table}


\section{Sensitivity of the window count to calibration error}
The calibrated window edges carry the uncertainty of the calibration. Shifting the window by plus or minus the
leave-one-out root-mean-square error of the calibrated $\Eks$ (\qty{0.115}{eV}) changes the number of valid records
inside the calibrated Kohn--Sham window from 46 to between 34 and 58. In this window-shift test the count moves by about
$\pm 25\,\%$, so counts produced by the gap step should not be read to the unit.
The upper edge is a convention as well: without the window, 566 molecules of the cascade have a model efficiency of at least 1.86\,\%, the lowest value among the 37 survivors, and 31 of them reach 5\,\% or more (calibrated optical gaps \qtyrange{2.23}{2.58}{eV}, maximum 8.34\,\%).
% L-2026-10-18-12 (stat-hardening.txt), -24 (window-redundancy.txt)

\section{Molecules that pass the calibrated cascade}
Table~\ref{tab:si-survivors} lists the 37 survivors. Two entries and one group illustrate what the cascade does and does
not catch. Pentacene (ID 8671) passes, with a calibrated gap of \qty{2.02}{eV}. The anthraquinone dye 22936 passes; no
gate tests reactivity or photostability, so dyes of this kind are not excluded. Seven survivors have
$\Esone$ more than \qty{0.5}{eV} below $\Eopt$; the record check flags only the opposite disagreement and leaves them in.
Only 10 of the 37 lie within the chemical domain of the calibration set, as defined by the six-descriptor Mahalanobis
criterion of the Methods; the calibrated gaps of the other 27 carry an
error that the in-class leave-one-out error does not describe. The table, with the full SMILES, is also deposited as
\texttt{si\_survivors.csv}.
% L-2026-10-18-11 (22936, 7/37 downward disagreement), -17 (domain); table: make_si_tables.py
{\footnotesize
\begin{longtable}{rlp{4.1cm}rrrrrrr}
\caption{The 37 molecules that pass every gate of the calibrated cascade, ordered by Scharber efficiency (upper bound,
PCBM acceptor). $\Eks$: stored PubChemQC gap; $\Eopt$: calibrated optical gap; $\Esone$: sTDA-xTB singlet energy;
PCE$_{0.65}$ and PCE$_\mathrm{G}$: efficiency with $\FF = 0.65$ and with Green's fill factor; $d_\mathrm{M}$:
Mahalanobis distance to the calibration set, $^{*}$ marking the 27 molecules beyond the domain threshold
($d_\mathrm{M} > 4.29$). Energies in eV, $\Voc$ in V, efficiencies in \%.}\label{tab:si-survivors}\\
\toprule
ID & Formula & SMILES & $\Eks$ & $\Eopt$ & $\Esone$ & $\Voc$ & PCE$_{0.65}$ & PCE$_\mathrm{G}$ & $d_\mathrm{M}$\\
\midrule
\endfirsthead
\toprule
ID & Formula & SMILES & $\Eks$ & $\Eopt$ & $\Esone$ & $\Voc$ & PCE$_{0.65}$ & PCE$_\mathrm{G}$ & $d_\mathrm{M}$\\
\midrule
\endhead
17574 & C23H19N3O6S & \texttt{\seqsplit{CC(=O)N(c1ccc(cc1)Nc1cc(c(c2c1C(=O)c1ccccc1C2=O)N)S(=O)(=O)O)C}} & 2.33 & 1.98 & 1.66 & 1.60 & 10.1 & 14.3 & 4.3$^{*}$ \\
12919 & C13H10N4O4 & \texttt{\seqsplit{O=Nc1ccc(=O)c(cc1)N=NC(=c1ccn(cc1)O)O}} & 2.08 & 1.81 & 2.11 & 1.20 & 9.9 & 13.6 & 7.2$^{*}$ \\
12913 & C13H10N4O3 & \texttt{\seqsplit{O=Nc1ccc(c(=O)cc1)NNC(=O)c1ccncc1}} & 2.33 & 1.98 & 1.13 & 1.51 & 9.6 & 13.6 & 7.0$^{*}$ \\
18506 & C20H17N3O4 & \texttt{\seqsplit{CCCCn1c(=O)c2c(c1=O)c(N)c1c(c2N)c(=O)c2c(c1=O)cccc2}} & 2.26 & 1.93 & 1.50 & 1.39 & 9.5 & 13.3 & 3.0 \\
22936 & C22H17N3O6S & \texttt{\seqsplit{CC(=O)Nc1ccc(cc1)Nc1cc(c(c2c1C(=O)c1ccccc1C2=O)N)S(=O)(=O)O}} & 2.04 & 1.78 & 1.73 & 1.10 & 9.4 & 12.8 & 4.7$^{*}$ \\
327 & C14H8N2O8 & \texttt{\seqsplit{OC(=O)c1cc(C(=O)O)c2c([nH]1)c(O)c(=O)c1c2[nH]c(c1)C(=O)O}} & 2.42 & 2.04 & 1.65 & 1.59 & 9.2 & 13.0 & 8.7$^{*}$ \\
22916 & C20H14N2O5S & \texttt{\seqsplit{O=C1c2ccccc2C(=O)c2c1c(Nc1ccccc1)cc(c2N)S(=O)(=O)O}} & 2.29 & 1.95 & 1.70 & 1.39 & 9.2 & 12.8 & 4.2 \\
19385 & C8H7N3O2S & \texttt{\seqsplit{Nc1ccc(c(c1)O)C1N=NC(=S)O1}} & 2.51 & 2.10 & 1.52 & 1.68 & 8.8 & 12.5 & 7.0$^{*}$ \\
22902 & C23H20N2O5S & \texttt{\seqsplit{Cc1cc(C)c(c(c1)C)Nc1cc(c(c2c1C(=O)c1ccccc1C2=O)N)S(=O)(=O)O}} & 2.47 & 2.07 & 1.78 & 1.46 & 8.1 & 11.3 & 2.8 \\
6689 & C20H15N3O8S2 & \texttt{\seqsplit{Nc1c2C(=O)c3ccccc3C(=O)c2c(cc1S(=O)(=O)O)Nc1ccc(c(c1)S(=O)(=O)O)N}} & 2.37 & 2.01 & 1.69 & 1.31 & 8.0 & 11.1 & 6.0$^{*}$ \\
5773 & C19H17N3O4S2 & \texttt{\seqsplit{O=C(N[C@@H]1C(=O)N2[C@@H]1SCC(=C2C(=O)[O-])C[n+]1ccccc1)Cc1cccs1}} & 2.35 & 1.99 & 1.85 & 1.26 & 7.8 & 10.8 & 4.2 \\
6690 & C28H14N2O4 & \texttt{\seqsplit{O=c1c2c3[nH]c4ccc5c(c4[nH]c3ccc2c(=O)c2c1cccc2)c(=O)c1c(c5=O)cccc1}} & 2.26 & 1.93 & 2.36 & 1.09 & 7.5 & 10.3 & 3.2 \\
5570 & C7H7NO2 & \texttt{\seqsplit{C[n+]1cccc(c1)C(=O)[O-]}} & 2.56 & 2.13 & 1.94 & 1.20 & 6.0 & 8.3 & 4.9$^{*}$ \\
20549 & C18H18N2O4 & \texttt{\seqsplit{OCCNc1ccc(c2c1C(=O)c1c(C2=O)cccc1)NCCO}} & 2.53 & 2.11 & 1.55 & 1.15 & 6.0 & 8.2 & 3.8 \\
17157 & C7H7NO2 & \texttt{\seqsplit{[O-]C(=O)c1cc[n+](cc1)C}} & 2.55 & 2.13 & 1.93 & 1.18 & 6.0 & 8.2 & 4.9$^{*}$ \\
16356 & C9H8N2O & \texttt{\seqsplit{CNC1=c2ccccc2=NC1=O}} & 2.56 & 2.13 & 0.37 & 1.19 & 5.9 & 8.2 & 4.8$^{*}$ \\
18514 & C18H18N2O6 & \texttt{\seqsplit{OCCNc1ccc(c2c1C(=O)c1c(C2=O)c(O)ccc1O)NCCO}} & 2.44 & 2.05 & 1.86 & 1.03 & 5.8 & 7.9 & 4.7$^{*}$ \\
17191 & C17H16N2O3 & \texttt{\seqsplit{OCCNc1ccc(c2c1C(=O)c1ccccc1C2=O)NC}} & 2.54 & 2.12 & 1.54 & 1.11 & 5.6 & 7.7 & 3.4 \\
6817 & C13H10N2O & \texttt{\seqsplit{Cn1c2ccccc2nc2c1cccc2=O}} & 2.63 & 2.18 & 2.34 & 1.19 & 5.5 & 7.6 & 5.1$^{*}$ \\
6680 & C21H15NO3 & \texttt{\seqsplit{Cc1ccc(cc1)Nc1ccc(c2c1C(=O)c1ccccc1C2=O)O}} & 2.62 & 2.18 & 1.82 & 1.17 & 5.5 & 7.6 & 4.1 \\
11296 & C11H8O5 & \texttt{\seqsplit{Oc1c(=O)cc2c(c(c1O)O)c(O)ccc2}} & 2.34 & 1.98 & 1.85 & 0.80 & 5.1 & 6.7 & 5.6$^{*}$ \\
9406 & C32H18 & \texttt{\seqsplit{c1ccc2c(c1)cc1c(c2)cc2c(c1)cc1c(c2)c2cccc3c2c2c1cccc2cc3}} & 1.99 & 1.75 & 2.42 & 0.56 & 5.0 & 6.2 & 6.0$^{*}$ \\
2134 & C22H28N4O4 & \texttt{\seqsplit{OCCNCCNc1ccc(c2c1C(=O)c1c(C2=O)cccc1)NCCNCCO}} & 2.57 & 2.14 & 1.54 & 0.97 & 4.8 & 6.5 & 6.4$^{*}$ \\
10719 & C14H16 & \texttt{\seqsplit{CCc1ccc(c2c(c1)c(C)cc2)C}} & 2.20 & 1.89 & 1.45 & 0.66 & 4.8 & 6.2 & 8.8$^{*}$ \\
20419 & C20H14N2O2 & \texttt{\seqsplit{O=C1c2ccccc2C(=O)c2c1c(ccc2N)Nc1ccccc1}} & 2.56 & 2.13 & 1.74 & 0.93 & 4.7 & 6.3 & 4.6$^{*}$ \\
4212 & C22H28N4O6 & \texttt{\seqsplit{OCCNCCNc1ccc(c2c1C(=O)c1c(C2=O)c(O)ccc1O)NCCNCCO}} & 2.53 & 2.11 & 1.79 & 0.86 & 4.5 & 5.9 & 6.9$^{*}$ \\
9139 & C32H18 & \texttt{\seqsplit{c1ccc2c(c1)cc1c(c2)cc2c(c1)cc1c(c2)cc2c3c1ccc1c3c(cc2)ccc1}} & 2.18 & 1.88 & 2.21 & 0.60 & 4.4 & 5.7 & 6.0$^{*}$ \\
14650 & C15H12N2O2 & \texttt{\seqsplit{CNc1ccc(c2c1C(=O)c1ccccc1C2=O)N}} & 2.63 & 2.18 & 1.73 & 0.90 & 4.2 & 5.6 & 4.4$^{*}$ \\
9105 & C32H18 & \texttt{\seqsplit{c1ccc2c(c1)cc1c(c2)cc2c3c1ccc1c3c(cc2)c2c(c1)cc1c(c2)cccc1}} & 2.36 & 2.00 & 2.53 & 0.67 & 4.1 & 5.3 & 6.0$^{*}$ \\
8671 & C22H14 & \texttt{\seqsplit{c1ccc2c(c1)cc1c(c2)cc2c(c1)cc1c(c2)cccc1}} & 2.40 & 2.02 & 2.23 & 0.68 & 4.1 & 5.3 & 7.2$^{*}$ \\
10881 & C7H6O & \texttt{\seqsplit{O=c1cccccc1}} & 2.32 & 1.97 & 2.77 & 0.62 & 4.0 & 5.1 & 7.2$^{*}$ \\
17189 & C16H14N2O2 & \texttt{\seqsplit{CNc1ccc(c2c1C(=O)c1ccccc1C2=O)NC}} & 2.51 & 2.10 & 1.60 & 0.75 & 3.9 & 5.2 & 3.9 \\
19729 & C16H14N2O4 & \texttt{\seqsplit{CNc1ccc(c2c1C(=O)c1c(C2=O)c(NC)ccc1O)O}} & 2.53 & 2.11 & 2.05 & 0.75 & 3.8 & 5.0 & 4.3 \\
22937 & C26H18N2O6 & \texttt{\seqsplit{Oc1ccc(cc1)Nc1ccc(c2c1C(=O)c1c(C2=O)c(O)c(cc1)O)Nc1ccc(cc1)O}} & 2.45 & 2.06 & 1.72 & 0.50 & 2.8 & 3.5 & 4.6$^{*}$ \\
17190 & C14H12N4O2 & \texttt{\seqsplit{O=C1c2c(N)ccc(c2C(=O)c2c1c(N)ccc2N)N}} & 2.56 & 2.13 & 1.95 & 0.55 & 2.7 & 3.4 & 5.5$^{*}$ \\
19428 & C20H14 & \texttt{\seqsplit{C1CC=c2c(=C1)cc1c3c2cc2ccccc2c3C=C1}} & 2.46 & 2.07 & 2.24 & 0.48 & 2.7 & 3.3 & 7.6$^{*}$ \\
9175 & C22H16 & \texttt{\seqsplit{c1ccc2=CC3C(=c4cc5ccccc5cc4=CC3)C=c2c1}} & 1.74 & 1.57 & 1.44 & 0.17 & 1.9 & 1.7 & 7.5$^{*}$ \\
\bottomrule
\end{longtable}}


\section{Formal verification}
The three statements of the Methods were proved in Lean~4 (toolchain v4.34.1) with Mathlib at commit
\texttt{d13f23b723b8a846827a245b89c10fc7d3f11612} (tag v4.34.1). The build completes with no \texttt{sorry}, no
\texttt{admit} and no added axiom: \texttt{\#print axioms} reports, for each of the seven theorems below, only Lean's
standard axioms \texttt{propext}, \texttt{Classical.choice} and \texttt{Quot.sound}. An independent reader compared each
formal statement with its intended meaning and found all three faithful. Not formalised: repeated cut-offs in a cascade
(harmless, not stated), infinite families of conventions, Green's fill factor for $v < 0$, conditions under which the
robust set is non-empty (false in general, see F2), and every numerical result of the paper, which rests on computation
and is checked through the analysis record instead. The Lean project is part of the code deposit
(DOI 10.5281/zenodo.18201812).
% L-2026-10-18-10; gate re-run 2026-10-18: LEAN GATE PASS; #print axioms run 2026-10-18 (L4)

\paragraph{F1, order invariance.} \texttt{successiveFilter s ps} applies the pass/fail predicates \texttt{ps} to the
finite set \texttt{s} one after another. \texttt{successiveFilter\_eq\_filter} shows that the result is the set of
elements passing every predicate; \texttt{successiveFilter\_perm} shows that any reordering of the predicates gives the
same set. Order therefore changes only how many records each gate is credited with removing. The theorem does not
say whether a gate's observable is meaningful for a record; that is the role of the validity gate
(Section~\ref{sec:salts}).
\VerbatimInput[fontsize=\small,breaklines,frame=lines]{../../formal/Formal/F01_OrderInvariance.lean}

\paragraph{F2, robustness across conventions.} \texttt{Robust C Adm} is the set of molecules admitted under every
convention \texttt{c} of a finite family \texttt{C}. \texttt{mem\_Robust} and \texttt{Robust\_subset} state that it is
exactly the intersection of the admitted sets, and \texttt{Robust\_empty\_example} gives a two-convention family for
which it is empty: a molecule that survives every convention need not exist.
\VerbatimInput[fontsize=\small,breaklines,frame=lines]{../../formal/Formal/F02_Intersection.lean}

\paragraph{F3, Green's fill factor.} With $v$ the normalised open-circuit voltage,
$\FF(v) = \bigl(v - \ln(v + 0.72)\bigr)/(v + 1)$. \texttt{FF\_bounds} proves $0 < \FF(v) < 1$ for $v \ge 0$, and
\texttt{FF\_strictMonoOn} proves that $\FF$ increases strictly for $v \ge 1$. Negative $v$ is not covered.
\VerbatimInput[fontsize=\small,breaklines,frame=lines]{../../formal/Formal/F03_GreenFF.lean}

\end{document}
